% Source: https://artofproblemsolving.com/community/c3h2258725_sierpinski_triangle
% Stable identity: aops-topic-2258725-q01
The figures below are built as follows. Level $1$ is an equilateral triangle. To get level $n+1$ from level $n$, take each of the smallest upright triangles of level $n$ (there are $3^{n-1}$ of them) and join the midpoints of its sides, which splits it into three smaller upright triangles and one upside-down triangle in the middle.

\begin{asy}
usepackage("amsmath");
size(9cm);
defaultpen(fontsize(10pt));
void sier(pair a, pair b, pair c, int d){
  draw(a--b--c--cycle,linewidth(0.7));
  if(d>1){ pair ab=(a+b)/2, bc=(b+c)/2, ca=(c+a)/2;
    sier(a,ab,ca,d-1); sier(ab,b,bc,d-1); sier(ca,bc,c,d-1); }
}
for(int k=1;k<=3;++k){
  pair s=((k-1)*2.6,0);
  sier(s,s+(2,0),s+(1,sqrt(3)),k);
  label("Level "+string(k),s+(1,-0.1),S);
}
\end{asy}

A triangle \emph{in the figure} is any triangle, of any size and orientation, whose three sides lie along drawn segments. For example, level $1$ contains $1$ triangle and level $2$ contains $5$. How many triangles does level $n$ contain?
